\section{Onde o fluido se aplica}
\label{sec:aplicacoes}

O Gray-Scott é um sistema de reação-difusão, não um escoamento no sentido de Navier-Stokes. Ele
entra no conjunto de ``fluidos'' do The Well como um meio contínuo descrito por EDPs de
transporte (difusão) acopladas a fontes não lineares (reação). As aplicações do modelo e da
segmentação dos seus padrões são as seguintes.

\paragraph{Química} O modelo descreve a autocatálise cúbica $A + 2B \to 3B$, $B \to C$ num
reator alimentado continuamente \cite{grayscott1984}. Os padrões previstos por Pearson
\cite{pearson1993} foram observados experimentalmente no sistema ferrocianeto-iodato-sulfito em gel
\cite{lee1993}. Segmentar as manchas permite medir número, área, velocidade e divisão de
estruturas, que são as grandezas comparadas entre experimento e simulação.

\paragraph{Biologia e ecologia} Mecanismos de Turing explicam padrões de pigmentação (zebras,
peixes-zebra, conchas) e o espaçamento de estruturas no desenvolvimento \cite{kondo2010,murray2003},
assim como os padrões de vegetação em regiões semiáridas \cite{klausmeier1999}. Nesses domínios, a
máscara segmentada é a observável primária: fração ocupada, comprimento de onda e conectividade
dos padrões.

\paragraph{Computação gráfica: o pipeline gráfico} Reação-difusão é uma técnica clássica de
síntese procedural de texturas \cite{turk1991,witkin1991}. No pipeline gráfico, o sistema é
integrado no espaço de textura (ou diretamente na malha \cite{turk1991}), e o campo resultante é
mapeado em cor, rugosidade ou deslocamento na etapa de \emph{shading}. A segmentação de $B$ é o
que transforma o campo contínuo em regiões (por exemplo, manchas de pelagem com bordas nítidas e
máscaras de material). Ela também é usada para controlar a evolução do padrão quadro a quadro em
animações, o que corresponde ao filtro temporal Q10. Implementações em GPU (\emph{compute} ou
\emph{fragment shaders} com textura em \emph{ping-pong}) usam o mesmo esquema de diferenças com
contorno periódico (\emph{wrap}) adotado aqui.

\paragraph{Aprendizado de máquina científico} O The Well usa o Gray-Scott como benchmark de
previsão de campos (modelos substitutos) \cite{ohana2024well}. Uma segmentação confiável das
estruturas fornece métricas físicas interpretáveis para validar esses modelos, como a preservação
do número de manchas e a coerência temporal das frentes. Essas métricas complementam erros
pixel a pixel como o VRMSE.

\paragraph{Por que segmentar como imagem} Os campos do The Well estão em grades regulares, e o
protocolo da turma limita todos a no máximo 128 px no maior eixo. Assim, o mesmo conjunto de
arquiteturas de segmentação semântica de imagens pode ser aplicado a todos os sistemas físicos. As
três redes cobrem as três famílias dominantes: convolucional com \emph{skip connections} (U-Net),
\emph{Transformer} hierárquico com decodificador leve (SegFormer) e classificação de máscaras por
consultas (Mask2Former). A comparação entre sistemas físicos mede, portanto, como as propriedades de
cada campo (escala das estruturas, contraste e dinâmica) interagem com o viés indutivo de cada
família.
